\chapter*{Avant-propos}
\addcontentsline{toc}{chapter}{Avant-propos}

Ce document a deux objectifs.

Le premier est pédagogique : présenter, de manière aussi intuitive que rigoureuse, les deux grandes
familles de modèles génératifs par transport continu de distribution qui dominent aujourd'hui la
génération d'images — les \textbf{modèles de diffusion} (DDPM) et le \textbf{Flow Matching} — en
partant des bases (qu'est-ce qu'un modèle génératif, pourquoi la génération d'images est un problème
difficile) jusqu'aux détails d'implémentation (architecture U-Net, conditionnement par texte,
classifier-free guidance).

Le second objectif est documentaire : ce cours sert de référence théorique au dépôt de code qui
l'accompagne (\texttt{text-\linebreak[1]to-\linebreak[1]image-\linebreak[1]generative-\linebreak[1]models}),
une implémentation \emph{from scratch} en PyTorch des deux approches, entraînée sur le jeu de données
CelebA-Dialog (visages annotés de légendes en langage
naturel), conçue pour tourner entièrement en local sur un Mac Apple Silicon (accélération MPS). Le
chapitre~\ref{chap:project} fait explicitement le lien entre chaque brique théorique et le fichier de
code qui l'implémente.

\begin{intuition}
Le fil conducteur de tout ce document tient en une phrase : \textbf{un modèle génératif de ce type
n'apprend pas à dessiner une image directement} — il apprend une règle de transformation simple
(« comment transformer un peu de bruit en quelque chose d'un peu plus net ») et on obtient une image
en appliquant cette règle un grand nombre de fois. La diffusion et le flow matching ne diffèrent que
par la nature de cette règle et la façon dont on l'apprend.
\end{intuition}

\paragraph{Comment lire ce document.} Les chapitres~\ref{chap:generative} à~\ref{chap:architecture}
constituent le cours proprement dit et sont indépendants du projet : ils se lisent comme une
introduction générale aux modèles de diffusion et de flow matching. Le chapitre~\ref{chap:project} est
spécifique à l'implémentation et suppose les chapitres précédents acquis.
